\chapter*{Glosario de Wolfram Mathematica}
\addcontentsline{toc}{chapter}{Glosario de Wolfram Mathematica}
\chaptermark{Glosario de Wolfram Mathematica}


%\newcommand{\MatCmd}[1]{\NoCaseChange{\MatCmd{#1}}}
\DeclareRobustCommand{\MatCmd}[1]{\texttt{#1}}
\AddToNoCaseChangeList{\MatCmd}

Las prácticas de la asignatura se plantean como una extensión del trabajo realizado en las clases de teoría y problemas. El objetivo no es aprender un lenguaje de programación ni sustituir los razonamientos matemáticos por una salida de ordenador. Se utilizará \textit{Wolfram Mathematica} para \emph{comprobar cálculos}, \emph{explorar ejemplos}, \emph{representar gráficamente resultados} y abordar, una vez comprendido el procedimiento, casos cuyo cálculo manual resulte innecesariamente largo.

Por este motivo, el número de herramientas que se empleará durante el curso es deliberadamente reducido. Este glosario contiene las órdenes y la sintaxis necesarias para realizar los ejercicios marcados con \M a lo largo de los apuntes. No pretende ser un manual de Mathematica.

\section*{Sintaxis básica}

Mathematica distingue entre mayúsculas y minúsculas. Los nombres de las funciones incorporadas comienzan por mayúscula y sus argumentos se escriben entre corchetes. Así, por ejemplo,
\begin{mathematicacode}
Sin[x]
Sqrt[x]
Exp[x]
\end{mathematicacode}
representan, respectivamente, $\sin x$, $\sqrt{x}$ y $e^x$.

Las listas se escriben entre llaves. Esta misma notación permite introducir vectores y matrices:
\begin{mathematicacode}
u = {1, -2, 3}
A = {{1, 2}, {3, 4}}
\end{mathematicacode}
Por tanto, una matriz se introduce como una lista cuyas componentes son sus filas.

La multiplicación puede escribirse mediante un espacio o mediante \MatCmd{*}; las potencias se indican con \MatCmd{\string^}. Por ejemplo,
\begin{mathematicacode}
3 x^2
3*x^2
\end{mathematicacode}
son expresiones equivalentes.

Para plantear ecuaciones se utiliza \MatCmd{==}, no \MatCmd{=}. El signo \MatCmd{=} se reserva para asignar un valor a un símbolo:
\begin{mathematicacode}
a = 3
Solve[x^2 == a, x]
\end{mathematicacode}

Cuando resulte cómodo definir una función utilizaremos el patrón \MatCmd{x\_} y la asignación diferida \MatCmd{:=}:
\begin{mathematicacode}
f[x_] := x^3 - 3 x
f[2]
\end{mathematicacode}
No será necesario utilizar mecanismos de programación más avanzados en esta asignatura.

Algunas constantes y símbolos incorporados que aparecerán con frecuencia son
\begin{center}
\begin{tabular}{ll}
\toprule
\textbf{Mathematica} & \textbf{Significado}\\
\midrule
\MatCmd{Pi} & $\pi$\\
\MatCmd{I} & unidad imaginaria $i$\\
\MatCmd{Infinity} & $\infty$\\
\MatCmd{Sqrt[x]} & $\sqrt{x}$\\
\MatCmd{Exp[x]} & $e^x$\\
\MatCmd{Log[x]} & $\ln x$\\
\MatCmd{Sin[x]}, \MatCmd{Cos[x]}, \MatCmd{Tan[x]} & funciones trigonométricas\\
\MatCmd{ArcCos[x]} & $\arccos x$\\
\bottomrule
\end{tabular}
\end{center}
Los argumentos de las funciones trigonométricas se interpretan en \textbf{radianes}.

\section*{Herramientas de uso general}

\begin{description}

\item[\MatCmd{N}] Obtiene una aproximación decimal de una expresión exacta.
\begin{mathematicacode}
N[Pi/7]
N[Sqrt[2], 12]
\end{mathematicacode}
En el segundo ejemplo se solicita una aproximación con mayor precisión.

\item[\MatCmd{Simplify}] Simplifica una expresión algebraica. Es especialmente útil para comprobar que dos resultados escritos de forma diferente son equivalentes.
\begin{mathematicacode}
Simplify[(x^2 - 1)/(x - 1)]
Simplify[(x + 1)^2 - (x^2 + 2 x + 1)]
\end{mathematicacode}
La simplificación automática no sustituye el estudio del dominio: dos expresiones pueden coincidir donde ambas están definidas y, sin embargo, tener dominios diferentes.

\item[\MatCmd{Solve}] Resuelve simbólicamente ecuaciones o sistemas de ecuaciones.
\begin{mathematicacode}
Solve[x^2 - 5 x + 6 == 0, x]
Solve[{x + y == 3, 2 x - y == 0}, {x, y}]
\end{mathematicacode}
Las ecuaciones se escriben con \MatCmd{==} y un sistema se introduce como una lista de ecuaciones.

\item[\MatCmd{Plot}] Representa una o varias funciones reales en un intervalo.
\begin{mathematicacode}
Plot[x Exp[-x], {x, 0, 8}]
Plot[{Sin[x], Cos[x]}, {x, -2 Pi, 2 Pi}]
\end{mathematicacode}
Una gráfica debe utilizarse para explorar o comprobar el estudio analítico, no para sustituirlo. En particular, la apariencia de una función puede cambiar sustancialmente al modificar el intervalo de representación.

\end{description}

\section*{Números complejos}

\begin{description}

\item[\MatCmd{Re}, \MatCmd{Im}] Devuelven las partes real e imaginaria de un número complejo.
\begin{mathematicacode}
z = 3 - 4 I
Re[z]
Im[z]
\end{mathematicacode}

\item[\MatCmd{Abs}] Para un número complejo devuelve su módulo; para un número real coincide con el valor absoluto.
\begin{mathematicacode}
Abs[3 - 4 I]
\end{mathematicacode}

\item[\MatCmd{Arg}] Devuelve el argumento principal de un número complejo no nulo.
\begin{mathematicacode}
Arg[-1 + Sqrt[3] I]
N[Arg[-1 + Sqrt[3] I]]
\end{mathematicacode}

\end{description}

Para ecuaciones en números complejos se utilizará el mismo comando \MatCmd{Solve} estudiado anteriormente:
\begin{mathematicacode}
Solve[z^3 == -8 I, z]
\end{mathematicacode}

\section*{Funciones y cálculo}

\begin{description}

\item[\MatCmd{Limit}] Calcula simbólicamente un límite.
\begin{mathematicacode}
Limit[(x^2 - 1)/(x - 1), x -> 1]
Limit[(5 x^2 + 3 x - 1)/(-2 x^2 + 7), x -> Infinity]
\end{mathematicacode}
Para los límites laterales puede añadirse la opción \MatCmd{Direction} cuando sea necesaria. En las prácticas ordinarias se priorizarán los límites para los que baste la sintaxis anterior.

\item[\MatCmd{D}] Calcula derivadas simbólicas. El segundo argumento indica la variable respecto de la que se deriva.
\begin{mathematicacode}
D[x^2 Exp[-x], x]
D[Sin[x^2], x]
\end{mathematicacode}
Las derivadas de orden superior pueden obtenerse indicando el orden:
\begin{mathematicacode}
D[Exp[x] Sin[x], {x, 3}]
\end{mathematicacode}
Si se ha definido previamente una función, Mathematica también entiende la notación con primas:
\begin{mathematicacode}
f[x_] := x^3 - 3 x
f'[2]
f''[2]
\end{mathematicacode}

\item[\MatCmd{Series}] Obtiene un desarrollo de Taylor alrededor de un punto hasta el orden indicado.
\begin{mathematicacode}
Series[Sin[x], {x, 0, 5}]
Series[Exp[x], {x, 0, 4}]
\end{mathematicacode}
La salida incluye un término que representa el orden del resto.

\item[\MatCmd{Normal}] Elimina de una serie el término que representa el resto y deja únicamente el polinomio de Taylor. Se utilizará solo cuando sea necesario manipular o representar dicho polinomio.
\begin{mathematicacode}
p = Normal[Series[Exp[x], {x, 0, 4}]]
Plot[{Exp[x], p}, {x, -2, 2}]
\end{mathematicacode}

\item[\MatCmd{Integrate}] Calcula integrales simbólicas, tanto indefinidas como definidas.
\begin{mathematicacode}
Integrate[x Exp[-x], x]
Integrate[x Exp[-x], {x, 0, 4}]
\end{mathematicacode}
En una integral indefinida Mathematica proporciona una primitiva y, por tanto, no añade una constante de integración. Al escribir la solución matemática debe recordarse la constante $C$.

\item[\MatCmd{NIntegrate}] Aproxima numéricamente una integral definida. Se utilizará cuando interese un valor numérico o cuando la primitiva no pueda expresarse mediante funciones elementales sencillas.
\begin{mathematicacode}
NIntegrate[Exp[-x^2], {x, 0, 1}]
\end{mathematicacode}

\end{description}

\section*{Matrices, espacios vectoriales y aplicaciones lineales}

Los vectores y matrices se introducen mediante listas. El operador punto \MatCmd{.} representa el producto matricial y, para dos vectores, el producto escalar:
\begin{mathematicacode}
A = {{1, 2}, {3, 4}}
x = {2, -1}
A.x
x.x
\end{mathematicacode}

\begin{description}

\item[\MatCmd{Det}] Calcula el determinante de una matriz cuadrada.
\begin{mathematicacode}
Det[{{1, 2}, {3, 4}}]
\end{mathematicacode}

\item[\MatCmd{Inverse}] Calcula la inversa de una matriz regular.
\begin{mathematicacode}
Inverse[{{1, 2}, {3, 4}}]
\end{mathematicacode}

\item[\MatCmd{MatrixRank}] Calcula el rango de una matriz.
\begin{mathematicacode}
MatrixRank[{{1, 2, 3}, {2, 4, 6}, {0, 1, 1}}]
\end{mathematicacode}

\item[\MatCmd{RowReduce}] Obtiene la forma escalonada reducida por filas. Es especialmente útil para comprobar una eliminación de Gauss y para estudiar sistemas mediante su matriz ampliada.
\begin{mathematicacode}
RowReduce[{{1, 1, 3}, {2, -1, 0}}]
\end{mathematicacode}
En este ejemplo la última columna puede interpretarse como la columna de términos independientes de un sistema de dos ecuaciones.

\item[\MatCmd{LinearSolve}] Resuelve directamente un sistema matricial $Ax=b$ cuando se proporcionan la matriz de coeficientes y el vector de términos independientes.
\begin{mathematicacode}
A = {{1, 1}, {2, -1}}
b = {3, 0}
LinearSolve[A, b]
\end{mathematicacode}
Para discutir sistemas con parámetros o para obtener una descripción simbólica más flexible se utilizará normalmente \MatCmd{Solve}.

\item[\MatCmd{NullSpace}] Devuelve una base del núcleo de una matriz. Si $A$ representa una aplicación lineal, permite comprobar una base de $\ker A$.
\begin{mathematicacode}
A = {{1, 2, 3}, {2, 4, 6}}
NullSpace[A]
\end{mathematicacode}

\item[\MatCmd{Orthogonalize}] Aplica un proceso de ortonormalización a una familia de vectores linealmente independientes. Se utilizará para comprobar el procedimiento manual de Gram--Schmidt.
\begin{mathematicacode}
Orthogonalize[{{1, 1, 0}, {1, 0, 1}}]
\end{mathematicacode}

\item[\MatCmd{Transpose}] Calcula la matriz traspuesta.
\begin{mathematicacode}
Transpose[{{1, 2, 3}, {4, 5, 6}}]
\end{mathematicacode}
Será útil, entre otros lugares, para comprobar cambios por congruencia del tipo $P^t A P$.

\end{description}

\section*{Diagonalización}

\begin{description}

\item[\MatCmd{Eigenvalues}] Calcula los valores propios de una matriz cuadrada, contando multiplicidades algebraicas.
\begin{mathematicacode}
A = {{2, 1}, {1, 2}}
Eigenvalues[A]
\end{mathematicacode}

\item[\MatCmd{Eigenvectors}] Calcula vectores propios asociados a los valores propios de una matriz.
\begin{mathematicacode}
Eigenvectors[A]
\end{mathematicacode}
Estos comandos se utilizarán para \emph{comprobar} los valores y vectores propios calculados a mano. La decisión de si una matriz es diagonalizable debe justificarse mediante las multiplicidades y las dimensiones de los subespacios propios.

\end{description}

No se introducirá un comando específico que diagonalice automáticamente la matriz. Una vez construidas manualmente las matrices $P$ y $D$, la comprobación se realiza con las herramientas ya conocidas:
\begin{mathematicacode}
Inverse[P].A.P
P.D.Inverse[P]
\end{mathematicacode}
Para formas cuadráticas y diagonalización por congruencia se procede de forma análoga:
\begin{mathematicacode}
Transpose[P].A.P
\end{mathematicacode}

\section*{Geometría vectorial}

No utilizaremos comandos específicos de geometría analítica. Los problemas se resolverán con las mismas operaciones vectoriales estudiadas en el curso.

\begin{description}

\item[\MatCmd{.}] El punto entre dos vectores representa su producto escalar. Entre matrices, o entre una matriz y un vector, representa el producto matricial.
\begin{mathematicacode}
u = {1, 2, -1}
v = {2, 0, 3}
u.v
\end{mathematicacode}

\item[\MatCmd{Norm}] Calcula la norma euclídea de un vector.
\begin{mathematicacode}
Norm[{1, 2, -1}]
\end{mathematicacode}
De esta forma, el ángulo entre dos vectores no nulos puede obtenerse reutilizando el producto escalar:
\begin{mathematicacode}
ArcCos[(u.v)/(Norm[u] Norm[v])]
\end{mathematicacode}

\item[\MatCmd{Cross}] Calcula el producto vectorial de dos vectores de $\R^3$.
\begin{mathematicacode}
Cross[{1, 0, 2}, {0, 1, -1}]
\end{mathematicacode}

\end{description}

Las pertenencias, intersecciones y posiciones relativas que conduzcan a sistemas de ecuaciones se comprobarán con \MatCmd{Solve}, sin introducir órdenes adicionales de geometría.

\section*{Resumen de herramientas}

La siguiente tabla recoge el conjunto de órdenes que se utilizará de forma habitual. Una misma herramienta aparece en varios temas porque precisamente se pretende reutilizar un vocabulario computacional pequeño.

\begin{center}
\small
\begin{tabularx}{\textwidth}{>{\ttfamily}l X}
\toprule
\textbf{Herramienta} & \textbf{Uso principal}\\
\midrule
N & aproximaciones numéricas\\
Simplify & simplificación y comprobación de expresiones\\
Solve & ecuaciones y sistemas simbólicos\\
Plot & representación de funciones\\
Re, Im, Abs, Arg & números complejos\\
Limit & límites\\
D & derivadas\\
Series, Normal & polinomios de Taylor\\
Integrate & integrales simbólicas\\
NIntegrate & integrales numéricas\\
Det, Inverse & determinante e inversa\\
MatrixRank & rango\\
RowReduce & eliminación de Gauss\\
LinearSolve & sistemas lineales en forma matricial\\
NullSpace & núcleo\\
Orthogonalize & Gram--Schmidt\\
Transpose & traspuesta y cambios por congruencia\\
Eigenvalues, Eigenvectors & valores y vectores propios\\
. & producto escalar y producto matricial\\
Norm, Cross & norma y producto vectorial\\
\bottomrule
\end{tabularx}
\end{center}

Esta lista es intencionadamente corta. Siempre que un ejercicio pueda resolverse combinando estas herramientas se evitará introducir un nuevo comando que automatice directamente el procedimiento matemático que se pretende aprender.
